\section{Conclusion and future work}
\label{sec:conclusion}

We replicated the PCA statistical arbitrage strategy of \citet{avellaneda2010} on S\&P~500
stocks from 2011. The residual mean-reversion signal persists, with a gross in-sample Sharpe
ratio of \BaseGrossSR{}, but at 5~bps one-way cost the original specification loses money
because its gross edge is only \BaseBreakeven{}~bps per unit traded. About half of its trading
comes from hedging, and hedging fewer factors together with a no-trade band roughly doubles the
edge per trade, bringing the strategy to about break-even in-sample. That conclusion does not
survive realistic timing. Trading one day after the signal removes up to half of the gross Sharpe
ratio, and every rule-based variant then loses money. A deep model trained on the net Sharpe
ratio is much more robust to the delay and has the highest unselected break-even cost (about
\DeepLagBE~bps), but it also falls short of 5~bps, and part of its edge is a systematic tilt
rather than arbitrage. In-sample, the strategy family is profitable only if trading costs are
below about 4~bps.

Planned next steps:
\begin{enumerate}
  \item \textbf{Statistical evaluation:} deflated Sharpe ratio, block-bootstrap confidence
        intervals, cost-sensitivity curves, and analysis of performance across regimes.
  \item \textbf{Costs and execution:} per-stock spread estimates from daily high and low prices,
        and execution at the next open.
  \item \textbf{Extensions:} choosing the number of factors each day from the Marchenko--Pastur
        edge, shrinkage estimators of the correlation matrix \citep{ledoit2004}, and a deep model
        constrained to a balanced residual book.
  \item \textbf{Data:} survivorship-free CRSP data with delisting returns.
  \item \textbf{Final evaluation:} a single, pre-registered evaluation on the 2019--2025 held-out
        period.
\end{enumerate}
